\section{TRAPS Construction}
\label{sec:construction}

This section gives the TRAPS construction for one active epoch and the region-allocation objective. Layering trade-offs, key rotation, and migration details follow the epoch-update discipline below.

\subsection{Epoch Policy}

Fix an active epoch \(e\). The epoch has a trapdoor key \(K_e^{\mathrm{td}}\), a replay-suppression key \(K_e^{\mathrm{neg}}\), and a runtime policy
\[
\Pi_e
=
\left(
\phi_e,
\mathcal{O}_e,
\{\tau_j\}_{j=0}^{r},
\{\mathcal I_{e,j}\}_{j=1}^{r},
\varepsilon_{\mathrm{cap}}
\right).
\]
Here \(\phi_e\) is the stable feature extractor, \(\mathcal{O}_e\) is the learning and risk oracle, \(\{\tau_j\}\) are routing thresholds, \(\mathcal I_{e,j}\) is the region-local indexing substrate, and \(\varepsilon_{\mathrm{cap}}\) is the configured per-region residual forwarding cap.

For a credential attempt \((u,pw)\) with salt and metadata \((s,\sigma)\), TRAPS computes
\[
\begin{aligned}
x &= \phi_e(u,pw,\sigma), &
z &= \mathcal{O}_e(x), &
t &= f_{K_e^{\mathrm{td}}}(h(u,pw,s)).
\end{aligned}
\]
The score \(z\) selects region
\[
j=\pi_e(u,pw,\sigma)
\iff
\tau_{j-1}\le z<\tau_j.
\]
Routing depends on stable risk features; membership depends on the trapdoor. This separation keeps learning outside the authentication-correctness path.

\subsection{Credential Creation and Rotation}

When a password is registered or rotated, the service first applies its ordinary account-abuse and password-policy controls. Exact hard-deny checks, such as mandatory registration-time password-policy rejection, occur before TRAPS insertion. If the credential is accepted, TRAPS records the stable metadata needed by the active epoch, including any risk markers produced by known-password-list probes or chosen-password registration detectors.

TRAPS then computes the credential trapdoor
\[
t_u=f_{K_e^{\mathrm{td}}}(h(u,pw,s_u))
\]
and inserts \(t_u\) into the region selected by \(\pi_e(u,pw,\sigma_u)\). A risky but accepted credential is not treated as invalid; it is represented in a hardened region. This is essential for one-sided correctness: correct future logins using that credential must still be forwarded to the backend.

A credential becomes active only after the backend verifier, edge-directory entry, protected screening record, and region insertion have committed. This atomic activation rule prevents a valid credential from reaching the login path before it is represented.

\subsection{Login Screening}

For a login attempt \((u,pw)\), the frontend obtains \((\hat{s},\hat{\sigma},\hat v)\) from the edge directory. If \(u\) is unknown, it substitutes a dummy tuple \((\bar{s},\bar{\sigma},\bar v)\). TRAPS computes the region \(j\), trapdoor \(t\), and replay token
\[
\eta
=
\operatorname{HMAC}_{K_e^{\mathrm{neg}}}
\left(
e \parallel \hat v \parallel u \parallel t
\right).
\]

If the exact replay-suppression cache contains \(\eta\), the attempt is rejected locally with the generic failure response. Otherwise, TRAPS queries \(\mathcal I_{e,j}\). A negative result is locally rejected. A positive result is forwarded to the backend verifier \(H\). If the backend returns a definitive credential mismatch, TRAPS inserts \(\eta\) into the replay-suppression cache.

This path has one-sided screening correctness under the deployment invariant stated below: every active valid credential must be represented in at least one active queried epoch under the same routing decision used for insertion.

\subsection{Correctness and Adaptive Work Bound}
\label{sec:correctness-bound}

\paragraph{Operational contract}
The one-sided result is enforced by five deployment invariants. P1, stable route coverage: every represented credential is inserted into every epoch/region pair in \(\Gamma_e(u,pw,c)\), the complete online coverage set that the router may query for the credential. P2, directory completeness: the active directory maps every represented credential to at least one live region/filter version. P3, atomic activation: a directory/filter version becomes visible only after all represented credentials for that version have been inserted. P4, versioned migration: rebuild queries check all old/new versions in which the credential may reside, and old versions retire only after coverage is complete. P5, replay-cache safety: cache entries are exact backend-confirmed invalid tuples scoped by account and credential version, and are invalidated on account creation, password change, or credential-version change before the tuple can become represented.

\begin{theorem}[One-sided pre-screening]
Under P1--P5, TRAPS may forward or reject invalid tuples, but it never locally rejects a login attempt using a represented active credential. During maintenance or recovery, forwarding to the backend or querying a complete safe version preserves this contract until the next epoch is fully active.
\end{theorem}

\begin{proof}[Proof sketch]
For a represented credential, enrollment or migration inserted the same trapdoor into the same region that login recomputes from stable metadata and the submitted password. The exact replay cache cannot contain the current-version valid tuple because replay tokens are inserted only after backend mismatch. The region-local substrate is one-sided for inserted items, so it returns positive. If multiple epochs are active, TRAPS forwards when any active epoch returns positive.
\end{proof}

\begin{theorem}[Replay-capped invalid work]
For a fixed epoch with exact replay retention, if every region satisfies \(\varepsilon_j\le \varepsilon_{\mathrm{cap}}\), then an adaptive adversary issuing a set \(\mathcal U_A\) of distinct invalid credential tuples causes expected backend work
\[
\mathbb{E}[W_A]
\le
\sum_{x\in\mathcal U_A}\varepsilon_{j(x)}
\le
|\mathcal U_A|\varepsilon_{\mathrm{cap}} .
\]
Repeated submissions of a backend-rejected false positive do not add backend work during the replay-token retention window.
\end{theorem}

Here \(W_A\) counts backend verifier invocations induced by invalid tuples. The theorem therefore gives the direct per-distinct-tuple backend-work bound \(\mathbb E[W_A]/|\mathcal U_A|\le\varepsilon_{\mathrm{cap}}\), with replay discoveries consuming at most one backend verification while retained.

\subsection{Known-password and Chosen-password Routing}

TRAPS treats known-password lists and chosen-password detectors as risk signals rather than login-time validity decisions.

For known-password traffic, a risk-list probe can route attempts using leaked, weak, popular, or policy-denied passwords into a hardened region with a lower residual forwarding budget. If the probe misses some risky attempts, those attempts remain in a less hardened region; correctness is unaffected. If the probe falsely marks an attempt as risky, the attempt receives stricter screening or step-up treatment, but TRAPS does not accept or reject a credential solely because of that approximate signal.

For chosen-password pressure, registration-side detectors can mark accepted credentials as CPwA-suspect. These credentials are inserted into a hardened region, increasing its occupancy but preserving their validity. The allocation rule below reserves memory for this occupancy pressure and enforces a per-region cap so that steering into the hardened region remains bounded.

Risk-list updates that would change the routing of already active valid credentials are deployed through the epoch mechanism in this section. In particular, a newly introduced password-derived risk predicate cannot be used to reroute legacy credentials unless the old region remains queried or the credentials have been migrated.

\subsection{Risk-aware Region Allocation}

Let \(M\) be the total memory budget and \(M_j\) the memory assigned to region \(j\). For the Bloom-filter instantiation, define
\[
\beta_{\mathrm{BF}}=(\ln 2)^2.
\]
For provisioned occupancy \(\bar N_j\), the region residual forwarding probability is approximated by
\[
\varepsilon_j(M_j;\bar N_j)
\approx
\exp
\left(
-\beta_{\mathrm{BF}}\frac{M_j}{\bar N_j}
\right).
\]

The risk oracle induces attack masses over regions. Let \(q_j^{(\mathrm{K})}\) denote the fraction of distinct known-password invalid attempts routed to region \(j\), and let \(\alpha_j\) denote the estimated chosen-password insertion pressure in region \(j\). TRAPS uses the risk weight
\[
a_j
=
q_j^{(\mathrm{K})}
+
\lambda_{\mathrm{C}}\alpha_j,
\]
where \(\lambda_{\mathrm{C}}\ge 0\) controls how aggressively the system reserves memory for chosen-password occupancy pressure.

TRAPS chooses region memory by solving
\[
\begin{aligned}
\min_{\{M_j\ge 0\}}\quad
&\sum_{j=1}^{r} a_j\varepsilon_j(M_j;\bar N_j)\\
\text{s.t.}\quad
&\sum_{j=1}^{r}M_j\le M,\\
&\varepsilon_j(M_j;\bar N_j)\le \varepsilon_{\mathrm{cap}}
\quad \forall j .
\end{aligned}
\]
The objective minimizes weighted residual false forwards, while the cap ensures that no region becomes an arbitrarily weak steering target.

\begin{theorem}[Capped water-filling allocation]
For nonempty Bloom-backed regions with \(\bar N_j>0\), the allocation problem is feasible iff
\[
M\ge B_\Sigma
:=
\sum_j
\frac{\bar N_j}{\beta_{\mathrm{BF}}}
\ln\frac{1}{\varepsilon_{\mathrm{cap}}}.
\]
When feasible, an optimal solution assigns \(M_j=B_j+m_j\), where \(B_j=(\bar N_j/\beta_{\mathrm{BF}})\ln(1/\varepsilon_{\mathrm{cap}})\). Zero-weight regions receive \(m_j=0\). For \(a_j>0\),
\[
m_j
=
\frac{\bar N_j}{\beta_{\mathrm{BF}}}
\left[
\ln
\frac{a_j\varepsilon_{\mathrm{cap}}\beta_{\mathrm{BF}}}
{\mu\bar N_j}
\right]_+ ,
\]
with water level \(\mu>0\) chosen so that \(\sum_j m_j=M-B_\Sigma\).
\end{theorem}

\begin{proof}[Proof sketch]
The cap is equivalent to the lower bound \(M_j\ge B_j\), giving the feasibility condition. After assigning the base memory, the residual problem minimizes the separable convex function \(\sum_j a_j\varepsilon_{\mathrm{cap}}\exp(-\beta_{\mathrm{BF}}m_j/\bar N_j)\) over a simplex. The KKT conditions are sufficient. Positive residual allocations equalize marginal residual-forward reduction, and inactive coordinates are clipped at zero, yielding the stated water-filling rule.
\end{proof}

This is the allocation step evaluated as TRAPS learned allocation in Section~\ref{sec:experiments}.

The same control-layer objective can be instantiated with exact keyed fingerprints by choosing region-specific tag lengths \(b_j\), with residual collision budget approximately \(2^{-b_j}\). The empirical implementation in this paper uses the Bloom-filter instantiation.

\subsection{Dynamic Capacity Adaptation}
\label{sec:dynamic-capacity}

The static allocation above applies to one epoch. In deployment, the attack mixture can change: KPwA traffic can shift across password populations, CPwA registrations can increase occupancy in selected regions, and replay-cache observations can reveal where residual false forwards occur. TRAPS therefore treats region capacities as epoch targets produced by a control-plane optimizer.

Let \(w\) index a control window. TRAPS maintains smoothed estimates of distinct invalid traffic, CPwA-suspect insertions, and backend-confirmed false forwards per region. It converts these observations into a protected occupancy target
\[
\bar N_{j,w}^{+}
=
N_{j,w}
+
s_{j,w}
+
\chi_{\mathrm C}\overline\Delta_{j,w},
\]
and a dynamic risk weight
\[
a_{j,w}
=
\omega_{\mathrm K,w}\widehat q_{j,w}^{(\mathrm K)}
+
\omega_{\mathrm C,w}\widehat q_{j,w}^{(\mathrm C)}
+
\lambda_{\mathrm C}\widehat\alpha_{j,w}
+
\lambda_{\mathrm R}\widehat r_{j,w}.
\]
Here \(N_{j,w}\) is current represented occupancy, \(s_{j,w}\) is a safety margin, \(\overline\Delta_{j,w}\) is smoothed CPwA-suspect insertion pressure, \(\widehat q_{j,w}^{(\mathrm K)}\) and \(\widehat q_{j,w}^{(\mathrm C)}\) are estimated invalid-query mixtures, \(\widehat\alpha_{j,w}\) is normalized insertion pressure, and \(\widehat r_{j,w}\) is the normalized replay-cache false-forward discovery rate.

The next-epoch target solves the capped allocation problem with \(\bar N_{j,w}^{+}\) and \(a_{j,w}\). Optionally, the optimizer adds a convex transition penalty to avoid unnecessary capacity churn. Capacity reductions are realized only through epoch rebuild and dual-query migration, because shrinking an active Bloom filter in place could violate one-sided correctness. If the cap is infeasible under the current memory budget and projected occupancy, the control plane must increase memory, slow further CPwA-suspect insertions, relax the cap in degraded mode, or combine these actions.

The main replay figures use frozen capacities selected from training/validation statistics. Section~\ref{sec:exp-index-microbench} separately evaluates a next-epoch dynamic-capacity and drift-guard policy using previous-window region observations only; the dynamic and fallback rows use no future labels.

\subsection{Epoch Updates}

TRAPS distinguishes routing-preserving updates from routing-changing updates. Routing-preserving changes, such as changing layer splits inside a fully populated region, can be activated after the affected structures have been populated.

Routing-changing updates are deployed as new epochs. If the new routing decision can be computed from protected screening records, the new epoch is rebuilt off path and dual-queried during migration. If the new routing decision requires password-derived information that was not recorded at enrollment, the old epoch remains active and the credential migrates lazily after successful authentication, password rotation, or account recovery.

This update discipline preserves the central invariant: every active valid credential is represented in at least one active queried epoch under the routing rule used for insertion.
